\documentclass[a4paper,landscape,headrule,footrule]{foils}

%%
%%% macros for Theories of Grammar
%%%
\usepackage{polyglossia}
\setdefaultlanguage{english}
%\setmainfont{TeX Gyre Pagella}


\newcommand{\logo}{~}
\newcommand{\header}[3]{%
\title{\vspace*{-2ex} \large JMORF/MORS --- Morpho-Syntax
\\[2ex] \Large  \emp{#2} \\ \emp{#3}}
\author{\blu{Francis Bond}   \\ 
\normalsize  \textbf{Palacký University}\\
\normalsize  \url{https://fcbond.github.io/}\\
\normalsize  \texttt{bond@ieee.org}}
\MyLogo{JMORF (2026)}
\renewcommand{\logo}{#2}
\hypersetup{
   pdfinfo={
     Author={Francis Bond},
     Title={#1: #2},
     Subject={JMORF/MORS: Morpho-Syntax},
     Keywords={Syntax, Morphology, HPSG, Unification, Constructions},
     License={CC BY 4.0}
   }
 }

\date{#1
  \\ Location: KB-1.11}
}
\usepackage{hyperref}
\hypersetup{
     colorlinks,
     linkcolor={blue!50!black},
     citecolor={red!50!black},
     urlcolor={blue!80!black}
   }



\usepackage{xcolor}
\usepackage{graphicx}
\newcommand{\blu}[1]{\textcolor{blue}{#1}}
\newcommand{\grn}[1]{\textcolor{green}{#1}}
\newcommand{\hide}[1]{\textcolor{white}{#1}}
\newcommand{\emp}[1]{\textcolor{red}{#1}}
\newcommand{\txx}[1]{\textbf{\textcolor{blue}{#1}}}
\newcommand{\lex}[1]{\textbf{\mtcitestyle{#1}}}


% \renewcommand{\labelitemi}{\textcolor{violet}{‣}}
% \renewcommand{\labelitemii}{\textcolor{purple}{‣}}

\newcommand{\subhead}[1]{\noindent\textbf{#1}\\[5mm]}

\newcommand{\Bad}{\emp{\raisebox{0.15ex}{\ensuremath{\mathbf{\otimes}}}}}
\newcommand{\bad}[1]{*\eng{#1}}

\newcommand{\com}[1]{\hfill (\emp{#1})}%

\usepackage{relsize,xspace}
\newcommand{\into}{\ensuremath{\rightarrow}\xspace}
\newcommand{\tot}{\ensuremath{\leftrightarrow}\xspace}
\usepackage{url}
\newcommand{\lurl}[1]{\MyLogo{\url{#1}}}

\usepackage{langsci-gb4e}
\let\eachwordone=\slshape
%%% langsci-gb4e defines \thexnumi as \@roman\@xsi{xnumi}, which only works
%%% inside an exe environment; some decks use xlisti on its own.  Restore
%%% the plain gb4e definition so that keeps working.
\makeatletter
\def\thexnumi{\@xsi{xnumi}}
\makeatother
\newcommand{\lx}[1]{\textbf{\mtciteform{#1}}}
\newcommand{\ix}{\ex\slshape}



\newcommand{\ent}{\ensuremath{\Rightarrow}\xspace}
\newcommand{\ngv}{\ensuremath{\not\Rightarrow}\xspace}
%\usepackage{times}
%\usepackage{nttfoilhead}
\newcommand{\myslide}[1]{\foilhead[-25mm]{\raisebox{12mm}[0mm]{\emp{#1}}}\MyLogo{\logo}}
\newcommand{\myslider}[1]{\rotatefoilhead[-25mm]{\raisebox{12mm}[0mm]{\emp{#1}}}}
%\newcommand{\myslider}[1]{\rotatefoilhead{\raisebox{-8mm}{\emp{#1}}}}

\newcommand{\section}[1]{\myslide{}{\begin{center}\Huge \emp{#1}\end{center}}}



\usepackage[lyons,j,e,k]{mtg2e}
%\renewcommand{\mtcitestyle}[1]{\textcolor{-red!75!green!50}{\textsl{#1}}}
\renewcommand{\mtcitestyle}[1]{\textcolor{teal}{\textsl{#1}}}
\newcommand{\cs}{\mtciteform}
\newcommand{\iz}[1]{\texttt{\textup{#1}}}
\newcommand{\gm}{\textsc}
\usepackage[normalem]{ulem}
\newcommand{\ul}{\uline}
\newcommand{\ull}{\uuline}
\newcommand{\wl}{\uwave}
\newcommand{\vs}{\ensuremath{\Leftrightarrow}~}
%%%
%%% Bibliography
%%%
\usepackage{natbib}
%\usepackage{url}
\usepackage{bibentry}



\usepackage[utf8]{inputenc}
%%% trees
\usepackage{multicol}
%%% pst-node: one example in lec-12-control draws a connector between
%%% two points with \rnode/\ncbar.  This used to come in with rtrees.
\usepackage{pst-node}
\usepackage{forest}
\useforestlibrary{linguistics}
\forestapplylibrarydefaults{linguistics}

%%% \makebox's optional arguments are read as AVM delimiters by
%%% langsci-avm, so hide them behind a macro with no visible brackets.
\newcommand{\lbox}[1]{\makebox[0mm][l]{#1}}

%%% Tree styles.
%%%   words  -- trees whose leaves are actual words (italicised, aligned)
%%%   cat    -- trees whose leaves are categories or feature structures
\forestset{
  words/.style={for tree={s sep=7mm, l sep=5mm,
                          if n children=0{font=\itshape, tier=terminal}{}}},
  cat/.style={for tree={s sep=7mm, l sep=5mm}},
}
\newcommand{\wtree}[1]{\begin{forest} words #1 \end{forest}}
\newcommand{\ctree}[1]{\begin{forest} cat #1 \end{forest}}
%%% the rule being applied at this step of a derivation
\newcommand{\step}[1]{\begin{center}\emp{#1}\end{center}}

%%% AVMs (Language Science Press)
\usepackage{langsci-avm}

%%% Macros the decks use, carried over from headx.tex
\newcommand{\blank}{\rule{3em}{1pt}\xspace}
\newcommand{\ft}[1]{\textsc{#1}}
\newcommand{\idn}{\mtciteform}    %%% Indonesian
\newcommand{\prd}[1]{\textbf{#1}}
\newcommand{\typ}[1]{\textit{#1}}
%%% \val: provide before renewing, in case a package defines it.
\providecommand{\val}[1]{}
\renewcommand{\val}[1]{\textit{#1}}
%%% Manning's \avmbox, still used in tree labels and in maths where a
%%% bare langsci \tag would not work.  \avm{\tag{...}} gives the same
%%% boxed tag as the ones inside feature structures.
\newcommand{\avmbox}[1]{\mbox{\avm{\tag{#1}}}}

%%% Ring a value in a dotted oval.  These used pstricks \ovalnode,
%%% which came in with rtrees; tikz (via forest) does the same job.
\newcommand{\OV}[1]{\tikz[baseline]{\node[draw,dotted,red,ellipse,
          inner sep=1pt,outer sep=0pt]{#1};}}
\newcommand{\OVB}[1]{\tikz[baseline]{\node[draw,dotted,blue,ellipse,
          inner sep=1pt,outer sep=0pt]{#1};}}

%%% From CSLI book
\newcommand{\mc}{\multicolumn}
\newcommand{\HD}{\textbf{H}\xspace}
\newcommand{\el}{< >}       %%% empty list; langsci-avm needs the space

\newcommand{\A}{\noindent\textbf{A}: }
\newcommand{\Q}{\noindent\textbf{Q}: }

%%% Question and answer on two pages, written once:
%%%   \qaslide{Title}{... question ... \answer{... answer ...}}
%%% prints the slide twice, first without the \answer parts and then
%%% with them, so the answer appears "on click".  The example counter
%%% is reset so both pages show the same example numbers.
\newif\ifshowanswer
\newcommand{\answer}[1]{\ifshowanswer#1\fi}
\newcounter{qaxnumi}
\newcommand{\qaslide}[2]{%
  \setcounter{qaxnumi}{\value{xnumi}}%
  \showanswerfalse\myslide{#1}#2%
  \setcounter{xnumi}{\value{qaxnumi}}%
  \showanswertrue\myslide{#1}#2%
  \showanswerfalse}
%\newcommand{\C}{\noindent\textbf{C}: }



\begin{document}

\header{Lecture 7}{The Structure of the Lexicon I}{Lexical Types}
\maketitle



\myslide{Overview}
\begin{enumerate}
\item \txx{Lexical Types}: why a lexical hierarchy, the seven lexical
  terms, defaults, a tour of the lexeme hierarchy, the Case Constraint
  \\ Activity: all the constraints on \lex{give}
\item \txx{Lexical Rules: A First Look}: how lexical rules fit into the
  grammar, the three kinds of rule, and two examples
\end{enumerate}
\bigskip
Next week: lexical rules in detail, and the passive. \\
No problem set this week: the mid-term is due instead.


%%%%%%%%%%%%%%%%%%%%%%%%%%%%%%%%%%%%%%%%%%%%%%%%%%%%%%%%%%%%%%%%%%%%%%%%
%%% Converted from Emily Bender's slides (ch08a-uw.pdf), with permission,
%%% and updated to follow her 2026 deck (1026_lexicaltypes.pdf); the
%%% %%% pN comments give the page in that deck.
%%%%%%%%%%%%%%%%%%%%%%%%%%%%%%%%%%%%%%%%%%%%%%%%%%%%%%%%%%%%%%%%%%%%%%%%
\section{Lexical Types}

%%% p4
\myslide{Overview}
\begin{itemize}
\item Motivation for lexical hierarchy
\item Lexical terminology
\item Defaults: the `\,/\,' notation
\item Tour of the lexeme hierarchy
\item The Case Constraint
\item \val{pos} vs.\ \val{lexeme}
\item In-class activity: all the constraints on \lex{give}
\end{itemize}

%%% p6
\qaslide{Motivation for the Lexical Hierarchy}{
\begin{itemize}
\item We've streamlined our grammar rules\ldots
  \begin{itemize}
  \item \ldots by stating some constraints as general principles
  \item \ldots and locating lots of information in the lexicon.
  \end{itemize}
\item Our lexical entries currently stipulate a lot of
  information that is common across many entries and
  should be stated only once.  Examples?
  \answer{%
  \begin{itemize}
  \item \emp{Every common noun: \avm{[ head & [\type*{noun} agr & [ per & 3rd ] ] ]},
      \avm{[ spr & < DP > ]}, \avm{[ mode & ref ]}}
  \item \emp{Every verb: \avm{[ head & \type{verb} ]}, \avm{[ mode & prop ]},
      an NP subject first on \ft{arg-st}}
  \item \emp{Every noun and verb: the SHAC}
  \end{itemize}}
\end{itemize}
}

\myslide{The Ideal Lexical Entry}
\begin{itemize}
\item Ideally, particular lexical entries need only give
  phonological form, the semantic contribution,
  and any constraints truly idiosyncratic to the
  lexical entry.
  \begin{itemize}
  \item We don't get quite that far in this book, because of our
    `cutesy' role names (\ft{lover}, \ft{giver}, \ldots): the entry
    still has to say which argument plays which role
  \item Implemented grammars such as the ERG use generic role labels
    (\ft{arg1}, \ft{arg2}), so the mapping from arguments to roles can
    live in the lexical type
  \end{itemize}
\end{itemize}

%%% p7
\myslide{Lexemes and Words (in a General Linguistic Sense)}
\begin{itemize}
\item \txx{Lexeme}: An abstract proto-word which gives rise
  to genuine words. We refer to lexemes by their
  `dictionary form', e.g.\ `the lexeme \eng{run}' or `the
  lexeme \eng{dog}'.
\item \txx{Word}: A particular pairing of form and meaning.
  \eng{Running} and \eng{ran} are different words
\end{itemize}

%%% p8
\qaslide{Lexical Entries, Lexical Types \& Lexical Rules}{
\begin{itemize}
\item \txx{Lexical entries} capture the similarities among \lex{run, runs,
    running,} and \lex{ran}.
\item The \txx{lexical type hierarchy} captures the similarities among
  \lex{run, sleep,} and \lex{laugh}, among those and other verbs like
  \lex{devour} and \lex{hand}, and among those and other words like
  \lex{book}.
  \begin{itemize}
  \item[Q] What do \lex{devour} and \lex{book} have in common?
  \answer{\item[A] \emp{The SHAC: each agrees with its specifier}}
  \end{itemize}
\item \txx{Lexical rules} capture the similarities among \lex{runs, sleeps,
    devours, hands,}\ldots
\end{itemize}
}

%%% The seven terms, with #1--#2 in bold and the rest in grey; all
%%% plain when the range is 1--7.
\makeatletter
\newcommand{\lexterm}[4]{%
  \ifnum#2=1 \ifnum#3=7 #4\else\lexterm@hl{#1}{#2}{#3}{#4}\fi
  \else\lexterm@hl{#1}{#2}{#3}{#4}\fi}
\newcommand{\lexterm@hl}[4]{%
  \ifnum#1<#2 \textcolor{gray}{#4}\else
  \ifnum#1>#3 \textcolor{gray}{#4}\else\textbf{#4}\fi\fi}
\makeatother
\newcommand{\lexterms}[2]{%
\begin{itemize}
\item \lexterm{1}{#1}{#2}{Lexeme}
\item \lexterm{2}{#1}{#2}{Lexical entry}
\item \lexterm{3}{#1}{#2}{Lexical sequence}
\item \lexterm{4}{#1}{#2}{`Family of lexical sequences'}
\item \lexterm{5}{#1}{#2}{Lexical rule}
\item \lexterm{6}{#1}{#2}{Lexical rule instantiation}
\item \lexterm{7}{#1}{#2}{Word structure}
\end{itemize}}

%%% p9
\myslide{Lexical Terminology}
The book uses seven related terms
\lexterms{1}{7}

%%% p10
\myslide{Lexical Terminology}
\lexterms{1}{4}

%%% p11
\myslide{\val{lexeme}}
\begin{itemize}
\item A \emp{type} in our grammar
\item The root of the sub-hierarchy of all the lexeme types we use in
  defining lexical entries
\item So `a lexeme' loosely refers to the type \val{lexeme} or one of
  its subtypes
\item These types are generalizations about combinations of
  constraints
  \begin{itemize}
  \item e.g.\ the lexeme type \val{pn-lxm} encodes the generalization
    that proper nouns are (usually) singular
  \item the lexeme type \val{stv-lxm} encodes the generalization that
    a transitive verb takes two arguments
  \end{itemize}
\end{itemize}

%%% p12
\myslide{Lexical Entry}
\begin{itemize}
\item A piece of the grammar that we write in the lexicon
  \begin{itemize}
  \item e.g.\ `look, we have \lex{book}, it's of type \val{cntn-lxm},
    and we have \lex{give}, it's of type \val{dtv-lxm}'
  \end{itemize}
\end{itemize}
\begin{center}
\scalebox{0.72}{%
\begin{tabular}{@{}l@{\qquad}l@{}}
\avm{< \textnormal{book} , [\type*{cntn-lxm}
    sem & [ index & $i$ \\
            restr & < [ reln & book \\ inst & $i$ ] > ] ] >}
&
\avm{< \textnormal{give} , [\type*{dtv-lxm}
    sem & [ index & $s$ \\
            restr & < [ reln & give \\ sit & $s$ \\
                        giver & $i$ \\ recipient & $j$ \\ gift & $k$ ] > ] \\
    arg-st & < X$_i$, Y$_j$, Z$_k$ > ] >}
\end{tabular}}
\end{center}

%%% p13
\myslide{Lexical Sequence}
\begin{itemize}
\item A \emp{fully resolved feature structure} that is a pair of some
  orthographic or phonological string with some feature structure of
  type \val{lexeme} or \val{word}
\item There is a type \val{l-sequence} to encode this
  \begin{itemize}
  \item the \ft{input} and \ft{output} features of lexical rules have
    values of type \val{l-sequence}
  \end{itemize}
\end{itemize}
\begin{center}
\begin{small}
\begin{tabular}{|l|l|l|}\hline
\textsc{type} & \textsc{features/constraints} & \textsc{ist} \\ \hline
\val{list} & & \val{feat-struc} \\ \hline
\val{list($\tau$)} & \avm{[ first & $\tau$ \\ rest & \type{list($\tau$)} ]}
  & \val{list} \\ \hline
\val{l-sequence} & \avm{[ first & \type{atom} \\
                         rest & < \type{word} > $|$ < \type{lexeme} > ]}
  & \val{list} \\ \hline
\end{tabular}
\end{small}
\end{center}
\emp{In plain English: an \val{l-sequence} is a two-element list,
  a string followed by one \val{word} or \val{lexeme}.}

%%% p14
\qaslide{Lexical Sequences?}{
\begin{center}
\scalebox{0.62}{%
\avm{< \textnormal{book} , [\type*{cntn-lxm}
  syn & [ head & [\type*{noun}
                  agr & \1 [\type*{3sing} per & 3rd \\ num & sg \\ gend & neut ] \\
                  case & nom ] \\
          val & [ spr & < [ agr & \1 ] > \\ comps & < > \\ mod & < > ] ] \\
  sem & [ mode & ref \\ index & $i$ \\
          restr & < [ reln & book \\ inst & $i$ ] > ] \\
  arg-st & < DP$_i$[ count & + ] > ] >}
\qquad
\avm{< \textnormal{book} , [\type*{cntn-lxm}
  syn & [ head & [\type*{noun}
                  agr & \1 [\type*{plural} per & 3rd \\ num & pl ] \\
                  case & acc ] \\
          val & [ spr & < [ agr & \1 ] > \\ comps & < > \\ mod & < > ] ] \\
  sem & [ mode & ref \\ index & $i$ \\
          restr & < [ reln & book \\ inst & $i$ ] > ] \\
  arg-st & < DP$_i$[ count & + ] > ] >}}
\end{center}
\begin{itemize}
\item[Q] Technically, neither of these is a lexical sequence.  Why?
\answer{\item[A] \emp{Lexical sequences must be \textbf{fully
      resolved}, and these don't fully resolve the DP on the
      \ft{arg-st} and \ft{spr} lists.}}
\end{itemize}
}

%%% p15
\myslide{`Family of Lexical Sequences'}
\begin{itemize}
\item A pairing of orthographic form with some \emp{underspecified}
  feature structure of type \val{lexeme} or \val{word}, which can be
  fully resolved to a number of different lexical sequences
\item The book often refers to lexical entries, for example, as being
  `families of lexical sequences', because they license infinitely
  many fully resolved lexical sequences
\item So the examples on the previous slide are technically all
  families of lexical sequences, because the \ft{arg-st} and \ft{val}
  constituents are not fully specified
\end{itemize}

%%% p16
\myslide{Lexical Terminology}
\lexterms{5}{7}

%%% p17
\myslide{Lexical Rule}
\begin{itemize}
\item A part of the grammar that we write that allows us to derive\ldots
  \begin{itemize}
  \item inflected forms of words, given lexical entries
  \item or additional lexical sequences where the feature structure is
    some subtype of \val{lexeme}
    \begin{itemize}
    \item \ldots which can unify with the \ft{input} feature of
      additional lexical rules until we arrive at a lexical sequence
      pairing a word-form with a feature structure of type \val{word}
    \end{itemize}
  \end{itemize}
\end{itemize}

\myslide{Lexical Rule: An Example}
\begin{center}
3rd-Singular Verb Lexical Rule

\medskip
\avm{[\type*{i-rule}
  input & < \3 , [\type*{verb-lxm} sem & [ restr & \tag{A} ] ] > \\
  output & < F$_{3SG}$(\3) ,
            [ syn & [ head & [ form & fin \\ agr & \type{3sing} ] ] \\
              sem & [ restr & \tag{A} $\oplus$ \ldots ] \\
              arg-st & < [ case & nom ], \ldots > ] > ]}
\end{center}
\emp{In plain English: take a verb lexeme, add the 3rd singular
  ending, and make it a finite word with a nominative subject that
  agrees with it.}

%%% p18
\myslide{Lexical Rule Instantiation}
\begin{itemize}
\item A \emp{fully resolved} feature structure that corresponds to
  one of those rules in the grammar
\item That is, the rule with a particular lexical sequence as its
  \ft{input} and the corresponding lexical sequence as its \ft{output}
  \begin{itemize}
  \item e.g.\ the 3rd-Singular Verb Lexical Rule with
    $\langle$\textnormal{sleep}, \ldots$\rangle$ as \ft{input} and
    $\langle$\textnormal{sleeps}, \ldots$\rangle$ as \ft{output}
  \item one rule in the grammar; one instantiation per word it licenses
  \end{itemize}
\end{itemize}

%%% p19
\myslide{Word Structure}
\begin{minipage}[c]{0.5\textwidth}
\begin{itemize}
\item The bottom part of the tree
\item A fully resolved feature structure of type \val{word} over some
  string which is the spelling or pronunciation of the word
\item e.g.\ \eng{park} in \eng{the cat sleeps in the park}
\end{itemize}
\end{minipage}%
\begin{minipage}[c]{0.45\textwidth}
\begin{center}
\scalebox{0.6}{%
\ctree{[{\avm{[\type*{word}
  syn & [ head & [\type*{noun}
                  agr & \1 [\type*{3sing} per & 3rd \\ num & sg \\ gend & neut ] \\
                  case & acc ] \\
          val & [ spr & < \2 > \\ comps & < > \\ mod & < > ] ] \\
  sem & [ mode & ref \\ index & $j$ \\
          restr & < [ reln & park \\ inst & $j$ ] > ] \\
  arg-st & < \2 > ]}}
  [{\eng{park}}]]}}
\end{center}
\end{minipage}


%%% p22--30 in Emily's deck: the full discussion of default
%%% inheritance is in lec-08b.  Here just enough to read the hierarchy.
\myslide{Defaults: the `\,/\,' Notation}
\begin{itemize}
\item Generalizations with exceptions are common:
  \begin{itemize}
  \item Most nouns in English aren't marked for \ft{case}, but
    pronouns are.
  \item Most proper nouns are singular, but some (\lex{the Alps},
    sports teams) are plural.
  \item Most verbs take an NP subject first on \ft{arg-st}, but a few
    (Chapter 11) do not.
  \end{itemize}
\item A constraint written with a slash, e.g.\ \avm{[ num & / sg ]},
  is a \txx{default}: it holds unless a subtype or a lexical entry
  says otherwise.
\item A constraint without a slash is \txx{inviolable}: a subtype
  that contradicts it can have no instances.
\item Defaults are all resolved in the lexicon: a \val{word} in a
  tree never has a `\,/\,' in it.
\item More on how defaults combine next week.
\end{itemize}

%%% The lexeme hierarchy is shown repeatedly, with an arrow pointing
%%% at the type under discussion.  #1: comma-separated names of nodes
%%% to colour red; #2: TikZ code drawn after the tree (the arrow).
\makeatletter
\newcommand{\lxhl}[2]{\ifcsname lexhl@#1\endcsname\textcolor{red}{#2}\else #2\fi}
\newcommand{\lexhier}[2]{%
\begingroup
\@for\lx@n:=#1\do{\expandafter\def\csname lexhl@\lx@n\endcsname{}}%
\begin{center}
\scalebox{0.7}{\begin{forest}
cat, for tree={font=\itshape, s sep=3mm}
[{\lxhl{synsem}{synsem}\\\textnormal{\avm{[ \textsc{syn, sem} ]}}}, align=center, name=synsem, s sep=4mm
  [{lexeme\\\textnormal{\avm{[ arg-st ]}}}, align=center, name=lexeme
    [{\lxhl{infl}{infl-lxm}}, name=infl
      [{\lxhl{verb}{verb-lxm}}, l*=3, name=verb
        [{siv-lxm}, name=siv]
        [{piv-lxm}, name=piv]
        [{tv-lxm}, name=tv
          [{stv-lxm}, name=stv]
          [{dtv-lxm}, name=dtv]
          [{ptv-lxm}, name=ptv]]]
      [{\lxhl{cn}{cn-lxm}}, l*=3, name=cn
        [{\lxhl{cntn}{cntn-lxm}}, name=cntn]
        [{\lxhl{massn}{massn-lxm}}, name=massn]]]
    [{const-lxm}, name=const
      [{adj-lxm}, l*=2, name=adj]
      [{conj-lxm}, l*=2, name=conj]
      [{det-lxm}, l*=2, name=det]
      [{predp-lxm}, l*=2, name=predp]
      [{argmkp-lxm}, l*=2, name=argmkp]
      [{\lxhl{pn}{pn-lxm}}, l*=1.4, name=pn]
      [{\lxhl{pron}{pron-lxm}}, l*=1.4, name=pron]]]
  [{expression}, name=expression, l*=0.8
    [{word\\\textnormal{\avm{[ arg-st ]}}}, align=center, l*=0.6, name=word]
    [{\lxhl{phrase}{phrase}}, l*=0.6, name=phrase]]]
#2
\end{forest}}
\end{center}
\endgroup}
\makeatother
%%% the orange arrow
\tikzset{hlarrow/.style={overlay, ->, >=stealth, orange, line width=2mm}}

%%% p31
\myslide{Our Lexeme Hierarchy}
\lexhier{}{}

%%% p33
\myslide{Functions of Types}
\begin{itemize}
\item Stating what features are appropriate for
  what categories
\item Stating generalizations
  \begin{itemize}
  \item Constraints that apply to (almost) all instances
  \item Generalizations about selection --- where
    instances of that type can appear
  \end{itemize}
\end{itemize}

%%% p34
\myslide{Every \val{synsem} has the Features SYN and SEM}
\lexhier{synsem}{\draw[hlarrow] ([xshift=-25mm]synsem.west) -- ([xshift=-2mm]synsem.west);}

%%% p35
\myslide{No ARG-ST on \val{phrase}}
\lexhier{phrase}{\draw[hlarrow] ([xshift=20mm,yshift=-12mm]phrase.south) -- ([yshift=-2mm]phrase.south);}
\vspace*{-2ex}
\emp{\ft{arg-st} is declared on \val{lexeme} and on \val{word}, so
  only lexical items have it; a phrase never does.}

%%% p36
\myslide{A Constraint on \val{infl-lxm}: the SHAC}
\lexhier{infl}{\draw[hlarrow] ([xshift=-25mm]infl.west) -- ([xshift=-2mm]infl.west);}

%%% p37
\myslide{A Constraint on \val{infl-lxm}: the SHAC}
\vspace*{2em}
\begin{center}
\val{infl-lxm}: \avm{[ syn & [ val & [ spr & < [ agr & \1 ] > ] \\
                               head & [ agr & \1 ] ] ]}
\end{center}
\bigskip
\emp{In plain English: every inflecting lexeme (nouns and verbs)
  has the same \ft{agr} value as its specifier.  Stated once here, it
  is inherited by every noun and verb in the lexicon.}

%%% p38
\myslide{Constraints on \val{cn-lxm}}
\lexhier{cn}{\draw[hlarrow] ([xshift=25mm]cn.east) -- ([xshift=2mm]cn.east);}

%%% p39
\myslide{Constraints on \val{cn-lxm}}
\begin{center}
\val{cn-lxm}: \avm{[ syn & [ head & [\type*{noun}
                                     agr & [ per & 3rd ] ] \\
                             val & [ spr & < [ head & \type{det} \\
                                               index & \textit{i} ] > ] ] \\
                     sem & [ mode & / ref \\
                             index & \textit{i} ] \\
                     arg-st & < X > $\oplus$ / < > ]}
\end{center}
\emp{In plain English: a common noun is 3rd person, takes a determiner
  as its specifier, shares its index with that determiner, is
  normally referential, and normally has no arguments beyond the
  specifier.}

%%% p40
\myslide{Formally Distinguishing Count vs.\ Mass Nouns}
\lexhier{cntn,massn}{%
  \draw[hlarrow] ($(cntn.south)!0.5!(massn.south)+(11mm,-12mm)$) -- ([yshift=-2mm]cntn.south);
  \draw[hlarrow] ($(cntn.south)!0.5!(massn.south)+(13mm,-12mm)$) -- ([yshift=-2mm]massn.south);}

%%% p41
\myslide{Formally Distinguishing Count vs.\ Mass Nouns}
\vspace*{2em}
\begin{center}
\begin{tabular}{r@{\ }l}
\val{cntn-lxm}: & \avm{[ syn & [ val & [ spr & < [ count & + ] > ] ] ]}\\[3ex]
\val{massn-lxm}: & \avm{[ syn & [ val & [ spr & < [ count & $-$ ] > ] ] ]}
\end{tabular}
\end{center}
\bigskip
\emp{The only difference is which determiners they take
  (\eng{a book}, \eng{much furniture})}

%%% p42
\myslide{Constraints on \val{verb-lxm}}
\lexhier{verb}{\draw[hlarrow] ([xshift=-15mm,yshift=20mm]verb.north west) -- (verb.north west);}

%%% p43
\myslide{Constraints on \val{verb-lxm}}
\vspace*{2em}
\begin{center}
\val{verb-lxm}: \avm{[ syn & [ head & \type{verb} ] \\
                       sem & [ mode & prop ] \\
                       arg-st & / < NP, \ldots > ]}
\end{center}
\bigskip
\emp{In plain English: a verb lexeme has \ft{head} \val{verb}, expresses a proposition, and
  by default its first argument (its subject) is an NP.  The default
  lets a few verbs take other subjects (\eng{That Kim left surprised
    me}).}

%%% p44
\myslide{Subtypes of \val{verb-lxm}}
\begin{center}
\scalebox{0.8}{\begin{forest}
cat, for tree={font=\itshape, s sep=3mm, l sep=3mm}
[{verb-lxm}
  [{siv-lxm}]
  [{piv-lxm}]
  [{tv-lxm}
    [{stv-lxm}]
    [{dtv-lxm}]
    [{ptv-lxm}]]]
\end{forest}}
\end{center}
\vspace*{-1ex}
\begin{itemize}
\item \val{verb-lxm}: \avm{[ arg-st & / < NP, \ldots > ]}
  \begin{itemize}
  \item \val{siv-lxm}: \avm{[ arg-st & < NP > ]} \com{strict intransitive: \eng{sleep}}
  \item \val{piv-lxm}: \avm{[ arg-st & < NP, PP > ]} \com{prepositional intransitive: \eng{rely}}
  \item \val{tv-lxm}: \avm{[ arg-st & < NP, NP, \ldots > ]} \com{transitive}
    \begin{itemize}
    \item \val{stv-lxm}: \avm{[ arg-st & < NP, NP > ]} \com{strict transitive: \eng{devour}}
    \item \val{dtv-lxm}: \avm{[ arg-st & < NP, NP, NP > ]} \com{ditransitive: \eng{hand}}
    \item \val{ptv-lxm}: \avm{[ arg-st & < NP, NP, PP > ]} \com{prepositional transitive: \eng{put}}
    \end{itemize}
  \end{itemize}
\end{itemize}

%%% p45
\myslide{Proper Nouns and Pronouns}
\lexhier{pn,pron}{%
  \draw[hlarrow] ([xshift=4mm,yshift=-18mm]pn.south) -- ([yshift=-2mm]pn.south);
  \draw[hlarrow] ([xshift=12mm,yshift=-18mm]pron.south) -- ([yshift=-2mm]pron.south);}

%%% p46
\myslide{Proper Nouns and Pronouns}
\begin{center}
\begin{tabular}{r@{\ }l}
\val{pn-lxm}: &
\avm{[ syn & [ head & [\type*{noun}
                       agr & [ per & 3rd \\
                               num & / sg ] ] ] \\
       sem & [ mode & ref ] \\
       arg-st & / < > ]}\\[4ex]
\val{pron-lxm}: &
\avm{[ syn & [ head & \type{noun} ] \\
       sem & [ mode & / ref ] \\
       arg-st & < > ]}
\end{tabular}
\end{center}
\emp{In plain English: proper nouns are 3rd person, usually singular
  (but \eng{the Alps}) and usually take no specifier (but \eng{the
    Kim I know}).  Pronouns are nouns with no arguments, usually
  referential (but \eng{itself} is \avm{[ mode & ana ]}).}

%%% p47
%%% a green tick (no amssymb \checkmark here)
\newcommand{\tick}{\tikz\draw[green!60!black, line width=2pt] (0,0.4em) -- (0.3em,0) -- (0.8em,0.9em);}
\qaslide{The Case Constraint}{
An outranked NP is \avm{[ case & acc ]}.

\medskip
Which of these NPs does it make accusative?
\begin{itemize}
\item object of verb \hfill \answer{\tick}\hspace*{3.2em}
\item second object of verb \hfill \answer{\tick}\hspace*{3.2em}
\item object of argument-marking preposition \hfill \answer{\tick}\hspace*{3.2em}
\item object of predicational preposition \hfill \answer{(\tick)}\hspace*{2.8em}
\end{itemize}
\answer{\par\medskip
  \emp{Everything after the first element of an
    \ft{arg-st} list is outranked.  The object of an argument-marking
    preposition has the same rank as its PP, so it is outranked by the
    verb's subject too; the object of a predicational preposition is
    outranked within the preposition's own \ft{arg-st}.}}
}

%%% p49
\qaslide{The Case Constraint, Continued}{
An outranked NP is \avm{[ case & acc ]}.
\begin{itemize}
\item Subjects of verbs
  \begin{itemize}
  \item Should we add a clause to cover nominative subjects, e.g.\
    un-outranked NPs are \avm{[ case & nom ]}?
    \answer{%
    \begin{itemize}
    \item \emp{No.}  \eng{We expect them to leave.} (Chapter 12)
      \\ \eng{them} is the \ft{spr} of \lex{leave} and a \ft{comps}
      member of \lex{expect}, so this would conflict
    \item Lexical rules for finite verbs will handle nominative subjects.
    \end{itemize}}
  \end{itemize}
\item Any other instances of case marking in English?
  \answer{\\ \emp{Possessives: \eng{my}, \eng{Kim's} (Chapter 6, Problem 4).}}
\item Does it apply to case systems in other languages?
  \answer{\\ \emp{No: the Case Constraint is an English-specific constraint.}}
\end{itemize}
}

%%% p50
\myslide{Apparent Redundancy of \val{pos} and \val{lexeme}}
\begin{itemize}
\item Why do we need both the \val{pos} subhierarchy
  and lexeme types?
\item \val{pos}:
  \begin{itemize}
  \item Applies to words and phrases; models the
    relationship between them
    \\ e.g.\ `verbs head VPs', `nouns head NPs'
  \item Constrains which features are appropriate
    (no \ft{aux} on \val{noun})
  \end{itemize}
\item \val{lexeme}:
  \begin{itemize}
  \item Generalizations about combinations of
    constraints
  \end{itemize}
\end{itemize}

%%% p51
\myslide{Lexical Types \& Lexical Rules}
\begin{itemize}
\item Lexical entries capture the similarities among \lex{run, runs,
    running,} and \lex{ran}.
\item The lexical type hierarchy captures the similarities among
  \lex{run, sleep,} and \lex{laugh}, among those and other verbs like
  \lex{devour} and \lex{hand}, and among those and other words like
  \lex{book}.
  \begin{itemize}
  \item What do \lex{devour} and \lex{book} have in common?  The SHAC
  \end{itemize}
\item Lexical rules capture the similarities among \lex{runs, sleeps,
    devours, hands,}\ldots
\end{itemize}

%%% p53--54: in-class activity
\qaslide{Activity: All the Constraints on \lex{give}}{
Here is the lexical entry for \lex{give}.  Which types does it
inherit from, and what does each one add?
\begin{center}
\scalebox{0.8}{\avm{< \textnormal{give} , [\type*{dtv-lxm}
    sem & [ index & $s$ \\
            restr & < [ reln & give \\ sit & $s$ \\
                        giver & $i$ \\ recipient & $j$ \\ gift & $k$ ] > ] \\
    arg-st & < X$_i$, Y$_j$, Z$_k$ > ] >}}
\end{center}
\answer{%
\begin{itemize}\setlength{\itemsep}{0pt}
\item \emp{\val{dtv-lxm} $\subset$ \val{tv-lxm}: \avm{[ arg-st & < NP, NP, NP > ]}}
\item \emp{\val{verb-lxm}: \avm{[ head & \type{verb} ]}, \avm{[ mode & prop ]},
    \avm{[ arg-st & / < NP, \ldots > ]}}
\item \emp{\val{infl-lxm}: the SHAC}
\item \emp{\val{lexeme}, \val{synsem}: the features \ft{arg-st}, \ft{syn}
    and \ft{sem} are appropriate}
\end{itemize}}
}

\myslide{Step 1: \val{dtv-lxm} and \val{tv-lxm}}
\begin{center}
\avm{< \textnormal{give} , [\type*{dtv-lxm}
    sem & [ index & $s$ \\
            restr & < [ reln & give \\ sit & $s$ \\
                        giver & $i$ \\ recipient & $j$ \\ gift & $k$ ] > ] \\
    arg-st & < \emp{NP}$_i$, \emp{NP}$_j$, \emp{NP}$_k$ > ] >}
\end{center}
\emp{\val{dtv-lxm} says the three arguments are all NPs: the three
  variables X, Y, Z are resolved to NP.}

\myslide{Step 2: \val{verb-lxm}}
\begin{center}
\avm{< \textnormal{give} , [\type*{dtv-lxm}
    syn & [ head & \emp{\type{verb}} ] \\
    sem & [ mode & \emp{prop} \\
            index & $s$ \\
            restr & < [ reln & give \\ sit & $s$ \\
                        giver & $i$ \\ recipient & $j$ \\ gift & $k$ ] > ] \\
    arg-st & < NP$_i$, NP$_j$, NP$_k$ > ] >}
\end{center}
\emp{\val{verb-lxm} adds \avm{[ head & \type{verb} ]} and
  \avm{[ mode & prop ]}.  Its default \avm{[ arg-st & / < NP, \ldots > ]}
  is already satisfied, so it adds nothing new.}

\myslide{Step 3: \val{infl-lxm} (the SHAC)}
\begin{center}
\avm{< \textnormal{give} , [\type*{dtv-lxm}
    syn & [ head & [\type*{verb} agr & \emp{\avm{\1}} ] \\
            val & [ spr & < [ agr & \emp{\avm{\1}} ] > ] ] \\
    sem & [ mode & prop \\
            index & $s$ \\
            restr & < [ reln & give \\ sit & $s$ \\
                        giver & $i$ \\ recipient & $j$ \\ gift & $k$ ] > ] \\
    arg-st & < NP$_i$, NP$_j$, NP$_k$ > ] >}
\end{center}
\emp{The SHAC identifies the verb's \ft{agr} with its specifier's.
  Nothing yet links \ft{spr} to \ft{arg-st}: the ARP is a constraint
  on \val{word}, so it applies only after a lexical rule.}

\myslide{Activity: All the Constraints on \lex{give}}
\begin{center}
\avm{< \textnormal{give} , [\type*{dtv-lxm}
    syn & [ head & [\type*{verb} agr & \1 ] \\
            val & [ spr & < [ agr & \1 ] > ] ] \\
    sem & [ mode & prop \\
            index & $s$ \\
            restr & < [ reln & give \\ sit & $s$ \\
                        giver & $i$ \\ recipient & $j$ \\ gift & $k$ ] > ] \\
    arg-st & < NP$_i$, NP$_j$, NP$_k$ > ] >}
\end{center}
\emp{Compare the entry we started from: \ft{head}, \ft{val},
  \ft{mode} and the NP types all came from the hierarchy.  The entry
  itself only had to say the form, the meaning and the type.}


%%%%%%%%%%%%%%%%%%%%%%%%%%%%%%%%%%%%%%%%%%%%%%%%%%%%%%%%%%%%%%%%%%%%%%%%
%%% Converted from Emily Bender's slides (ch08b-uw.pdf), with permission,
%%% and updated to follow her 2026 deck (1028_lexicalrules.pdf); the
%%% %%% pN comments give the page in that deck.  The detailed treatment
%%% is in lec-08b-lexrules.tex.
%%%%%%%%%%%%%%%%%%%%%%%%%%%%%%%%%%%%%%%%%%%%%%%%%%%%%%%%%%%%%%%%%%%%%%%%
\section{Lexical Rules: A First Look}

\myslide{Overview}
\begin{itemize}
\item How lexical rules fit in: tree structures, lexical sequences,
  lexical licensing
\item Three kinds of lexical rule: inflectional, derivational,
  post-inflectional
\item Two examples: the Plural Noun LR and the Agent Nominalization LR
\item Advice on writing lexical rules
\end{itemize}

%%% p11--13: the definitions from SWB Chapter 9, (37)--(39)
\myslide{How Lexical Rules Fit In: Well-Formed Tree Structure}
$\Phi$ is a well-formed tree structure according to grammar G
if and only if:
\begin{enumerate}
\item $\Phi$ is a tree structure,
\item the label of $\Phi$'s root node satisfies the constraint
  \avm{[ syn & [ head & [\type*{verb} form & fin ] \\
                 val & [ comps & < > \\ spr & < > ] ] ]} and
\item each local subtree within $\Phi$ is either phrasally licensed
  or lexically licensed.
\end{enumerate}
\emp{A sentence is a finite clause with nothing missing, where every
  branching is licensed by a rule and every leaf by the lexicon.}

\myslide{How Lexical Rules Fit In: Lexical Sequences}
$\langle\omega, \phi\rangle$ is a lexical sequence if and only if
$\omega$ is a phonological form (an atom), $\phi$ is a feature
structure, and either:
\begin{enumerate}
\item G contains some lexical entry $\langle d_1, d_2\rangle$
  such that $\omega$ satisfies $d_1$ and $\phi$ satisfies $d_2$, or
\item there is some lexical rule instantiation licensed by G
  (a feature structure of type \val{l-rule}) whose \ft{output}
  value is $\langle\omega, \phi\rangle$.
\end{enumerate}
\emp{A form paired with a feature structure is a lexical sequence if
  it comes straight from a lexical entry, or is the output of a
  lexical rule.  The second clause is recursive: the rule's input
  must itself be a lexical sequence, so rules can feed each other.}

\myslide{How Lexical Rules Fit In: Lexical Licensing}
\begin{minipage}[c]{0.7\textwidth}
A word structure of the form on the right is licensed if and only if:
\begin{enumerate}
\item $\langle\omega, \phi\rangle$ is a lexical sequence, where
  $\phi$ is of type \val{word},
\item (Case Constraint:) an outranked NP is [\ft{case} acc], and
\item $\phi$ satisfies the Binding Theory.
\end{enumerate}
\end{minipage}%
\begin{minipage}[c]{0.25\textwidth}
\begin{center}
\ctree{[{$\phi$} [{$\omega$}]]}
\end{center}
\end{minipage}

\medskip
\emp{Only lexical sequences of type \val{word} can appear in trees.
  Lexemes have to pass through at least one lexical rule to become
  words: that is what the inflectional rules are for.}

%%% p8
\myslide{Parsimony \& Plausibility}
\begin{itemize}
\item Lexical rules capture \textbf{productive}
  generalizations.
  \begin{itemize}
  \item When a new word enters the lexicon, it is assigned a lexical
    type, and we expect every lexical rule that applies to that type
    to apply to the new word too.
  \end{itemize}
\item There may be some `precompiling'
  going on as well.
  \begin{itemize}
  \item Psycholinguistically, very frequent words may just be stored
    ready-made rather than derived each time.
  \end{itemize}
\end{itemize}

%%% p9
\qaslide{Three Kinds of Lexical Rules (very English-specific!)}{
\begin{itemize}
\item \txx{Inflectional}: \val{lexeme} to \val{word}
  (takes us `over the boundary' from lexeme to word)
  \\[1ex] Examples?
  \answer{\\ \emp{Singular Noun LR, Plural Noun LR,
      3rd-Singular Verb LR, Non-3rd-Singular Verb LR, Past-Tense Verb
      LR, Base Form LR, Constant Lexeme LR}}
\item \txx{Derivational}: \val{lexeme} to \val{lexeme}
  \\[1ex] Examples?
  \answer{\\ \emp{Present Participle LR, Past Participle LR, Agent
      Nominalization LR}}
\item \txx{Post-Inflectional}: \val{word} to \val{word}
  \\ (Chapters 11, 13, 14)
\end{itemize}
}

%%% p14
\myslide{Three Subtypes of \textit{l-rule}}
\begin{center}
\begin{small}
  \begin{forest}
    cat, for tree={font=\itshape, l sep=3mm}
    [{l-rule} [{i-rule}] [{d-rule}] [{pi-rule}]]
  \end{forest}

  \medskip
  \val{l-rule} :
  \avm{[ input & \type{l-sequence} < X , [ sem & / \2 ] > \\
         output & \type{l-sequence} < Y , [ sem & / \2 ] > ]}

  \medskip
  \val{i-rule} :
  \avm{[ input & < X , [\type*{lexeme}
                         syn & \3 \\
                         arg-st & \tag{A} ] > \\
         output & < Y , [\type*{word}
                          syn & \3 \\
                          arg-st & \tag{A} ] > ]}
  \qquad
  \val{d-rule} :
  \avm{[ input & < X , [\type*{lexeme}
                         syn & / \3 ] > \\
         output & < Y , [\type*{lexeme}
                          syn & / \3 ] > ]}
\end{small}
\end{center}
\begin{small}
\emp{\val{l-rule}: the output keeps the input's meaning, by default.
  \val{i-rule}: the output is a \val{word} with the same \ft{syn} and
  \ft{arg-st}, no exceptions.
  \val{d-rule}: the output is another \val{lexeme} with the same
  \ft{syn}, by default.}
\end{small}

%%% p15
\myslide{Plural Noun LR}
\begin{center}
  \avm{[\type*{i-rule}
    input & < \1 , \type{cntn-lxm} > \\
    output & < F$_{NPL}$(\1) , [\type*{word}
                syn & [ head & [ agr & [ num & pl ] ] ] ] > ]}
\end{center}
\emp{Take any count noun lexeme \avm{\1}, add the plural ending to
  its form, and make it a word whose number is plural.  That is all
  the rule itself says; everything else comes from the types.}


\myslide{A Derivational Rule: Agent Nominalization}
\begin{center}
  \avm{[\type*{d-rule}
    input & < \2 , [\type*{verb-lxm} sem & [ index & $s$ ] \\
                                    arg-st & < X$_i$ , NP$_j$ > ] > \\
    output & < F$_{\textit{-er}}$(\2) ,
              [\type*{cntn-lxm}
                sem & [ index & $i$ ] \\
                arg-st & < Y ( , PP$_j$[ form & of ] ) > ] > ]}
\end{center}
\emp{Add \textit{-er} to a strictly transitive verb lexeme.  The
  result is a count noun lexeme whose index is the verb's subject
  index $i$: a \lex{driver} is the one who drives.  The verb's object
  $j$ can only appear as an \eng{of}-PP: \eng{a builder of bridges}.}

%%% p36
\myslide{Practicalities --- Writing Lexical Rules}
\begin{itemize}
\item Determine the type of the LR.
\item Determine the class of possible inputs.
\item Determine what should change.
  \begin{small}
  \begin{itemize}
  \item If \ft{input} and \ft{output} values are identified (by default or otherwise) and
    only \ft{output} value is mentioned, then\ldots
    \\ \emp{information is added.}  e.g.\ Plural Noun LR
    \\[1ex] (Lexical sequences incompatible with that value are not possible inputs)
  \item If \ft{input} and \ft{output} values are identified by default, but different values
    are given on the \ft{input} and \ft{output} of the rule, then\ldots
    \\ \emp{information is changed.}  e.g.\ Agent Nominalization LR
  \item If \ft{input} and \ft{output} values are identified by an inviolable constraint, but
    different values are given on the \ft{input} and \ft{output} of the rule, then\ldots
    \\ \emp{there is no well-formed input: the rule can never apply.}
  \end{itemize}
  \end{small}
\end{itemize}


\myslide{Summary}
\begin{itemize}
\item Lexical entries say only what is idiosyncratic; the
  \txx{lexeme hierarchy} supplies the rest through inheritance
\item A `\,/\,' marks a \txx{default} that a subtype may override
\item \txx{Lexical rules} relate families of lexical sequences:
  \val{i-rule}s make words from lexemes, \val{d-rule}s make new
  lexemes, \val{pi-rule}s make words from words
\item A rule states only what changes; the rule's type and the
  input's type supply everything else
\item Next week: how the inherited constraints fill in a lexical rule,
  constant lexemes, \ft{arg-st} and \ft{form}, and the passive
\end{itemize}

\myslide{Acknowledgments and References}

\begin{itemize}
\item Course design and slides borrow heavily from Emily Bender's course:
\textit{Linguistics 566: Introduction to Syntax for Computational Linguistics}
\\ \url{http://courses.washington.edu/ling566}
\item The slides on lexical types and lexical rules are adapted from
  Emily Bender's, with her permission; the feature structures follow
  \citet{sag:wasow:bender:2003}, \copyright{} 2003 CSLI Publications
\item This lecture complements \citet[Chapter 8]{sag:wasow:bender:2003}
\end{itemize}


%\input{lec-05-questions}

\bibliographystyle{aclnat}
\bibliography{abb,mtg,nlp,ling}
\end{document}

%%% Local Variables: 
%%% coding: utf-8
%%% mode: latex
%%% TeX-PDF-mode: t
%%% TeX-engine: xetex
%%% End: 
